\documentclass[11pt,fontset=fandol]{ctexart}
\usepackage[a4paper,margin=2.2cm]{geometry}
\usepackage{amsmath,amssymb,enumitem,xcolor,tcolorbox}
\tcbuselibrary{breakable}
\setlength{\parindent}{0pt}
\setlength{\parskip}{0.45em}
\newcommand{\R}{\mathbb R}
\newenvironment{exercise}[2]{
  \begin{tcolorbox}[colback=white,colframe=black!35,title={#1\quad\small #2},breakable]
}{\end{tcolorbox}}

\title{函数基础短逻辑题：对应练习}
\author{}
\date{}

\begin{document}
\maketitle

本练习的题目顺序已经打乱。每道题右上角标注了它对应的原题题号；这里只给题干，不给解答。

\begin{exercise}{练习1}{对应原题第32题}
已知 $\lg(2a-1)=y$，其中 $2a-1>0$，用 $y$ 表示 $a$。
\end{exercise}

\begin{exercise}{练习2}{对应原题第7题}
求函数
\[
y=2x+\sqrt{3-x}
\]
的值域。
\end{exercise}

\begin{exercise}{练习3}{对应原题第24题}
已知 $f,g$ 都在 $\R$ 上严格递增。判断 $f-g$ 是否一定严格递增；若不一定，给出反例。
\end{exercise}

\begin{exercise}{练习4}{对应原题第49题}
化简
\[
\log_3 7\cdot\log_7 3.
\]
\end{exercise}

\begin{exercise}{练习5}{对应原题第18题}
已知 $f(x)=\sqrt{x+2}$，分别求 $f(x^2)$ 与 $[f(x)]^2$ 的定义域，并说明二者为何不同。
\end{exercise}

\begin{exercise}{练习6}{对应原题第40题}
已知对任意 $x,y\in\R$，
\[
f(x+y)+f(x-y)=2f(x)+2f(y)+1.
\]
求 $f(0)$。
\end{exercise}

\begin{exercise}{练习7}{对应原题第29题}
解不等式
\[
\left(\frac{1}{3}\right)^{x^2+1}\le \left(\frac{1}{3}\right)^{4x}.
\]
\end{exercise}

\begin{exercise}{练习8}{对应原题第13题}
判断下列两个函数是否为同一个函数：
\[
f(x)=\frac{x^2-4}{x-2},\qquad g(x)=x+2.
\]
\end{exercise}

\begin{exercise}{练习9}{对应原题第43题}
已知 $g$ 在区间 $I$ 上严格递减，且 $g(I)\subseteq J$；$f$ 在 $J$ 上严格递减。判断 $f(g(x))$ 在 $I$ 上的单调性并证明。
\end{exercise}

\begin{exercise}{练习10}{对应原题第1题}
已知函数 $f$ 在 $\R$ 上严格递减，且
\[
f(3x+1)<f(x-5).
\]
求 $x$ 的取值范围。
\end{exercise}

\begin{exercise}{练习11}{对应原题第39题}
若存在 $x\in[-1,3]$，使
\[
a=x^2-2x+4,
\]
求实数 $a$ 的取值范围。
\end{exercise}

\begin{exercise}{练习12}{对应原题第35题}
分别求 $f(x)=x^2$ 在 $[0,+\infty)$ 和 $(-\infty,0]$ 上的反函数。
\end{exercise}

\begin{exercise}{练习13}{对应原题第27题}
连续函数 $f$ 满足 $f(0)f(3)>0$。能否断定 $f$ 在 $(0,3)$ 内没有零点？若不能，请举反例。
\end{exercise}

\begin{exercise}{练习14}{对应原题第41题}
已知对任意 $x,y\in\R$，
\[
f(x+y)=f(x)+f(y),
\]
且 $f(2)=6$。求 $f(x+2)$，并由此求 $f(6)$。
\end{exercise}

\begin{exercise}{练习15}{对应原题第17题}
求函数
\[
f(x)=\sqrt{\frac{(x+2)^2}{x+2}}
\]
的定义域。
\end{exercise}

\begin{exercise}{练习16}{对应原题第47题}
判断方程
\[
2^x=-(x^2+1)
\]
是否有实数解，并说明理由。
\end{exercise}

\begin{exercise}{练习17}{对应原题第22题}
已知 $f$ 在 $\R$ 上严格递增，解方程
\[
f(2x^2-3)=f(5x+1).
\]
\end{exercise}

\begin{exercise}{练习18}{对应原题第50题}
若函数
\[
f(x)=\log_{a+2}x
\]
在 $(0,+\infty)$ 上递增，求实数 $a$ 的取值范围。
\end{exercise}

\begin{exercise}{练习19}{对应原题第5题}
已知函数 $f$ 满足
\[
f(4+x)=f(4-x)\qquad(x\in\R).
\]
证明其图像关于直线 $x=4$ 对称，并求 $f(-2)$ 与哪个函数值相等。
\end{exercise}

\begin{exercise}{练习20}{对应原题第31题}
解方程
\[
\lg(x+1)+\lg(x-2)=\lg(2x-1).
\]
\end{exercise}

\begin{exercise}{练习21}{对应原题第45题}
已知 $g(x)=f(3-x)$ 在 $[0,2]$ 上递增，判断 $f$ 在 $[1,3]$ 上的单调性。
\end{exercise}

\begin{exercise}{练习22}{对应原题第14题}
判断
\[
f(x)=\sqrt{(x-2)^2},\qquad g(x)=x-2
\]
是否为同一个函数。
\end{exercise}

\begin{exercise}{练习23}{对应原题第38题}
若对任意 $x\in[-2,1]$，都有
\[
x^2+2x-a\le0,
\]
求实数 $a$ 的取值范围。
\end{exercise}

\begin{exercise}{练习24}{对应原题第20题}
已知对任意 $x\in\R$，
\[
f(x+3)=-f(x).
\]
证明 $6$ 是 $f$ 的一个周期。
\end{exercise}

\begin{exercise}{练习25}{对应原题第9题}
方程
\[
\log_3(3x^2)=\log_3(x+4)
\]
与方程 $3x^2=x+4$ 的解集是否相同？说明理由。
\end{exercise}

\begin{exercise}{练习26}{对应原题第46题}
求使
\[
\sqrt{x+2}+\sqrt{4-x}\ge a
\]
对任意 $x\in[-2,4]$ 恒成立的 $a$ 的取值范围。
\end{exercise}

\begin{exercise}{练习27}{对应原题第3题}
已知偶函数 $f$ 在 $[0,+\infty)$ 上严格递增，比较
\[
f(3a+1)\quad\text{与}\quad f(a-2)
\]
的大小，并把答案写成关于 $a$ 的分类。
\end{exercise}

\begin{exercise}{练习28}{对应原题第42题}
已知对任意实数 $x,y$，
\[
f(x+y)=f(x-y)+4y.
\]
证明 $f(t)-2t$ 是常数。
\end{exercise}

\begin{exercise}{练习29}{对应原题第34题}
函数
\[
f(x)=x^2-4x\qquad(x\ge2)
\]
是否存在反函数？若存在，求其定义域。
\end{exercise}

\begin{exercise}{练习30}{对应原题第16题}
设
\[
f(x)=\frac{2x-1}{x+2}\quad(x\ne-2).
\]
判断 $2$ 是否属于 $f$ 的值域，并求 $f$ 的值域。
\end{exercise}

\begin{exercise}{练习31}{对应原题第28题}
已知 $f$ 在 $[a,b]$ 上连续，且 $f(a)f(b)<0$。有人据此断言 $f$ 在 $(a,b)$ 内恰有一个零点。指出问题，并说还需要补充什么条件。
\end{exercise}

\begin{exercise}{练习32}{对应原题第10题}
解方程
\[
2^{x+2}+4^x=20.
\]
\end{exercise}

\begin{exercise}{练习33}{对应原题第48题}
证明方程
\[
3^x=5-x
\]
有且只有一个实数解。
\end{exercise}

\begin{exercise}{练习34}{对应原题第26题}
已知偶函数 $f$ 的定义域为 $\R$，方程 $f(x)=-1$ 有且只有三个实根，其中一个是 $0$。另外两个实根有什么关系？
\end{exercise}

\begin{exercise}{练习35}{对应原题第2题}
已知 $f$ 在 $[-1,4]$ 上严格递减。由
\[
f(3x-2)>f(x+1)
\]
能否直接得到 $3x-2<x+1$？若不能，还缺什么？
\end{exercise}

\begin{exercise}{练习36}{对应原题第37题}
已知 $f$ 的唯一零点为 $2$，求函数
\[
g(x)=f(x^2-4x+1)
\]
的零点。
\end{exercise}

\begin{exercise}{练习37}{对应原题第44题}
函数 $f$ 在 $[0,+\infty)$ 上递增。能否据此判断 $f(x^2-1)$ 在 $\R$ 上递增？说明理由。
\end{exercise}

\begin{exercise}{练习38}{对应原题第11题}
已知 $f$ 是严格单调函数，且
\[
f^{-1}(3a+1)=a-2.
\]
若 $f(x)=2x+5$，求 $a$。
\end{exercise}

\begin{exercise}{练习39}{对应原题第25题}
已知奇函数 $f$ 在 $[-3,3]$ 上有最大值 $7$，求其最小值。
\end{exercise}

\begin{exercise}{练习40}{对应原题第30题}
解不等式
\[
\log_{1/3}(x+2)>\log_{1/3}(5-x).
\]
\end{exercise}

\begin{exercise}{练习41}{对应原题第15题}
设
\[
f(x)=|x+2|,
\qquad
g(x)=\begin{cases}-x-2,&x<-2,\\x+2,&x\ge-2.\end{cases}
\]
判断 $f,g$ 是否为同一个函数。
\end{exercise}

\begin{exercise}{练习42}{对应原题第36题}
已知函数 $f$ 的唯一零点为 $-1$。求 $g(x)=f(x+4)$ 的零点。
\end{exercise}

\begin{exercise}{练习43}{对应原题第8题}
已知函数 $f$ 在 $[0,5]$ 上连续且严格递增，并且 $f(0)f(5)<0$。证明方程 $f(x)=0$ 在 $(0,5)$ 内有且只有一个实根。
\end{exercise}

\begin{exercise}{练习44}{对应原题第23题}
已知 $f,g$ 都在区间 $I$ 上递增，证明 $2f+3g$ 在 $I$ 上递增。
\end{exercise}

\begin{exercise}{练习45}{对应原题第6题}
已知函数 $f$ 的定义域为 $[-2,1)$，求
\[
g(x)=f(x^2+x-1)
\]
的定义域。
\end{exercise}

\begin{exercise}{练习46}{对应原题第19题}
已知 $f(x+4)=f(x)$ 对任意实数 $x$ 成立，且 $f(2)=5$，求 $f(-10)$。
\end{exercise}

\begin{exercise}{练习47}{对应原题第21题}
已知 $f$ 是偶函数，且 $f(x+3)=f(x)$。求证
\[
f\left(\frac32+x\right)=f\left(\frac32-x\right).
\]
\end{exercise}

\begin{exercise}{练习48}{对应原题第12题}
已知对任意 $x\in[-1,3]$，不等式
\[
a\ge 2x-x^2+1
\]
恒成立，求 $a$ 的取值范围。
\end{exercise}

\begin{exercise}{练习49}{对应原题第33题}
已知 $\log_aM=p,\log_aN=q$，其中 $a>0,a\ne1,M,N>0$。不用直接引用对数运算法则，求 $\log_a(MN)$。
\end{exercise}

\begin{exercise}{练习50}{对应原题第4题}
已知
\[
f(x)=x^7+ax^5+bx^3-6,
\qquad f(-1)=10,
\]
求 $f(1)$。
\end{exercise}

\end{document}
